\documentclass[book.tex]{subfiles}
\begin{document}

\chapter{Linear Least Squares Regression}

\label{chap:linear_regression}
\noindent\rule{11cm}{0.4pt} \\
\textbf{Prerequisites:} \autoref{chap:vectors}, \autoref{chap:linear_systems} \\
\textbf{Difficulty Level:} * \\
\noindent\rule{11cm}{0.4pt}
\vspace{5mm}

In this chapter we will learn various criteria for identifying the best fit of a straight line to data. We will learn how to compute the slope and intercept of a best-fit line with linear regression.

\section{Introduction to Data Analysis}

Data analysis is a process of inspecting and modeling data with the goal of discovering useful information, suggesting conclusions, and supporting decision-making. In engineering, we often try to generalize data from experience and make meaningful inferences. Typically our hypothesis space will span models varying greatly in complexity: some models will have many more free parameters or degrees of freedom than others. Under traditional approaches to data analysis, we adjust each model's parameters until it best represents the data. We assess models with a natural eye for complexity, balancing fit to the data with model complexity in subtle ways that will not inevitably prefer the most complex model. Instead we often seem to judge models using the \textit{Occam's razor}, which says that when faced with two opposing explanations for the same set of evidence, our mind naturally prefers the explanation that makes the fewest assumptions. In the context of analyzing trends in data, we choose the least complex hypothesis that fits the data well. In doing so, we avoid \textit{over-fitting} our data in order to support successful generalizations and predictions.

Two types of applications are generally encountered in Data Analysis:
\begin{enumerate}
	\item\textit{Trend analysis} represents the process of using the pattern of the data to make predictions. This can involve extrapolation beyond the limits of the observed data or interpolation within the range of the data.
	\item In \textit{hypothesis testing}, an existing mathematical model is compared with measured data. If the model coefficients are unknown, it may be necessary to determine values that best fit the observed data.
\end{enumerate}
There are two general approaches for recognizing patterns in data that are distinguished from each other on the basis of the amount of error associated with the data. 
\begin{enumerate}
\item When the data exhibit a significant degree of error, the strategy is to derive a single curve that represents the general trend of the data. Because any individual data point may be incorrect, we make no effort to intersect every point. Rather, the curve is designed to follow the pattern of the points taken as a group. One approach of this nature is called \textit{least-squares regression}. This avoids over-fitting the data and captures the general trend.
\item Second, where the data are known to be very precise, the basic approach is to fit a curve or a series of curves that pass directly through each of the points. This curve can now be used to estimate values between well-known discrete points and is termed as \textit{interpolation}.
\end{enumerate}

\section{Linear Least Squares Regression}
Where substantial error is associated with data, the best curve-fitting strategy is to derive
an approximating function that fits the shape or general trend of the data without necessarily
matching the individual points. For a data which is approximately linear, a \textit{line} would be the curve that fits these data points. One approach to do this is to visually inspect the plotted data and then sketch a \textquotedblleft best\textquotedblright \hspace{2pt} line through the points. However, unless the points define a perfect straight line, different analysts would draw different lines.

To remove this subjectivity, some criterion must be devised to establish a basis for the
fit. One way to do this is to derive a curve that minimizes the discrepancy between the data points and the curve. To do this, we must first quantify the discrepancy. The simplest example is fitting a straight line to a set of paired observations: ($x_1$,$y_1$),($x_2$,$y_2$),($x_3$,$y_3$)...($x_n$,$y_n$). The mathematical expression for the straight line is given by
\begin{align}
y = a_0 + a_1x + e
\end{align}

\begin{marginfigure}
	\centering
	\includegraphics[width = 2.4in]{figures/part1b/linear_least_squares/inadequatebestfit.PNG}
	\centering
	\caption{Examples of some criteria for \textit{best fit} that are inadequate for regression: (a) minimizes the sum of the residuals, (b) minimizes the sum of the absolute values of the residuals, and (c) minimizes the maximum error of any individual point.}
	\label{fig:1}
\end{marginfigure}

Here $a_0$ and $a_1$ are the coefficients representing the intercept and the slope, respectively, and $e$ is the error, or \textit{residual}, between the model and the observations, which can be represented by rearranging the above equation as
\begin{align}
e = y - a_0 - a_1x
\end{align}
Thus, the residual is the discrepancy between the true value of $y$ and the approximate value, $a_0$ + $a_1x$, predicted by the linear equation.
\subsection{Criteria for the \textquotedblleft Best\textquotedblright \hspace{2pt} Fit}

One strategy for fitting a best line through the data would be to minimize the sum of the
residual errors for all the available data. However, this is an inadequate criterion, as illustrated by Figure \ref{fig:1}(a), which depicts the fit of a straight line to two points. Obviously, the best fit is the line connecting the points. However, any straight line passing through the midpoint of the connecting line (except a perfectly vertical line) results in a minimum value of equal to zero because positive and negative errors cancel.
\begin{align}
\sum_{i=1}^{n}e_i = \sum_{i=1}^{n}(y_i - a_0 - a_1x_i)
\end{align}
One way to remove the effect of the signs might be to minimize the sum of the absolute
values of the discrepancies. Figure \ref{fig:1}(b) demonstrates why this criterion is also inadequate. For the four points shown, any straight line falling within the dashed lines will minimize the sum of the absolute values of the residuals. Thus, this criterion also does not yield a unique best fit.
\begin{align}
\sum_{i=1}^{n}|e_i| = \sum_{i=1}^{n}|y_i - a_0 - a_1x_i|
\end{align}
A third strategy for fitting a best line is the \textit{minimax} criterion. In this technique, the line is chosen that minimizes the maximum distance that an individual point falls from the line. As depicted in Figure \ref{fig:1}(c), this strategy is ill-suited for regression because it gives undue influence to an \textit{outlier} i.e., a single point with a large error.
\subsection{Least-Squares Fit of a Straight Line}
A strategy that overcomes the shortcomings of the aforementioned approaches is to minimize the sum of the squares of the residuals.
\begin{align}
S_r = \sum_{i=1}^{n}e_i^2 = \sum_{i=1}^{n}(y_i-a_0-a_1x_i)^2
\end{align}
This criterion, which is called \textit{least squares}, has a number of advantages, including that it yields a unique line for a given set of data.  Determining the slope and intercept for the best fit is calculated as follows.

To determine values for $a_0$ and $a_1$, $S_r$ is differentiated with respect to each unknown coefficient:
\begin{align}
\frac{\partial S_r}{\partial a_0}& = -2\sum(y_i-a_0-a_1x_i)\\
\frac{\partial S_r}{\partial a_1}&= -2\sum(y_i-a_0-a_1x_i)x_i
\end{align}
Setting these derivatives equal to zero will result in a minimum $S_r$. If this is done, the equations can be expressed as
\begin{align}
0 &= \sum y_i-\sum a_0- \sum a_1x_i\\
0 &= \sum x_iy_i-\sum a_0x_i- \sum a_1x_i^2
\end{align}

Now, realizing that $\sum a_0$ = $na_0$, we can express the equations as a set of two simultaneous linear equations with two unknowns ($a_0$ and $a_1$):
\begin{align}
na_0 + \sum x_i(a_1) &= \sum y_i\\
\sum x_i(a_0) + \sum x_i^2(a_1) &=\sum x_iy_i
\end{align}
These are called the \textit{normal equations}. They can be solved simultaneously for
\begin{align}
a_1 &= \dfrac{n\sum x_iy_i - \sum x_i\sum y_i}{n\sum x_i^2-(\sum x_i)^2}
\end{align}
\begin{align}
a_0 &= \bar{y} - a_1\bar{x}
\end{align}
where $\bar{y}$ and $\bar{x}$ are the means of $y$ and $x$, respectively.

\begin{example}
The housing prices with their area in square feet at Seattle, WA are given in a dataset. Determine the average price of the one square feet using Linear Least Squares Regression. Use the trend to make a bid at a house of 4000 square feet.
\textit{Matlab Code: } 
\begin{Verbatim}
load('Housing_Prices.mat')

Sum_x = sum(Area);
Sum_y = sum(Price);
Sum_xsq = sum(Area.*Area);
Sum_xy = sum(Area.*Price);

n = length(Area);
A = [0 (n*Sum_xsq - (Sum_x)^2);1 (Sum_x/n)];
B = [(n*Sum_xy - Sum_x*Sum_y) ;Sum_y/n];
a = A\B;
\end{Verbatim}

\begin{marginfigure}
	\centering
	\includegraphics[width = 2.7in]{figures/part1b/linear_least_squares/bestfit.png}
	\centering
	\caption{Data vs the Least Square Fit of the Data}
	\label{fig:2}
\end{marginfigure}

The slope obtained is $a_1$ = 280.6236 dollars (which is the average cost of one square feet) and the intercept is $a_0$ = -4.3581e+04. 

Using the linear trend, the bid to buy a house of 4000 square feet would be placed at 

$a_0$ + $a_1\times4000$ $=$ $ 1.079 $ million dollars 

The line along with the data can be observed from the Figure \ref{fig:2}.
%\end{document}
\end{example}
\section{Comments on Linear Regression}
It should be noted that the following statistical assumptions are inherent in linear least-squares procedures:
\begin{enumerate}
	\item Each $x$ has a fixed value; it is not random and is known without error.
	\item The $y$ values are independent random variables and all have the same variance.
	\item The $y$ values for a given $x$ must be normally distributed.
\end{enumerate}
Such assumptions are relevant to the proper derivation and use of regression. For
example, the first assumption means that (1) the $x$ values must be error-free and (2) the
regression of $y$ versus $x$ is not the same as $x$ versus $y$.




\end{document}